% !TeX spellcheck = ru_RU
\section{Заключение}

В ходе выполнения работы достигнуты следующие результаты:
\begin{itemize}
    \item Спроектирована и реализована модульная архитектура, в которой парсер конкретного фандома представлен классом, наследующим абстрактный \texttt{BaseParser}. Замена фандома осуществляется заменой одной строки в файле запуска без изменения DS-части.
    \item Разработаны два парсера: для вселенной «Властелин колец» (20 страниц) и для вселенной «Леди Баг» (15 страниц). Код парсера включает защиту от блокировки (User-Agent, паузы).
    \item Реализован модуль извлечения сущностей на базе spaCy, выделяющий персонажей (\texttt{PER}) и локации (\texttt{LOC}), присваивающий каждой уникальной сущности идентификатор и формирующий таблицы \texttt{entities.csv} и \texttt{doc\_entities.csv}. Проведён эксперимент с альтернативной библиотекой Natasha, однако она не дала улучшений.
    \item Данные из CSV загружены в PostgreSQL через DataGrip; созданы таблицы с автоинкрементом и внешними ключами. Проблемы импорта (ошибка «пустой идентификатор в кавычках») решены пересохранением файлов в UTF-8 с полным экранированием полей.
    \item Написаны SQL-запросы, отвечающие на контрольные вопросы для обеих вселенных: топ-10 персонажей локации, список локаций персонажа, часто упоминаемые вместе, документы по сущности.
    \item Построен граф совместных упоминаний персонажей для вселенной LOTR с фильтрацией (степень >= 5, вес ребра >= 5). После фильтрации граф содержит 23 вершины и 107 рёбер. Крупные узлы соответствуют центральным персонажам (Фродо, Гэндальф, Саурон, Арагорн).
\end{itemize}

Таким образом, поставленная цель достигнута. Система готова к использованию для любого другого фандома (например, Warhammer) путём создания нового класса парсера.

Исходный код проекта доступен в репозитории \url{https://github.com/s-vera666/fandom_wiki_parser}. В перспективе планируется добавить нормализацию падежных форм для улучшения SQL-запросов, а также автоматический сбор ссылок на все страницы фандома (вместо ручного прописывания). Это позволит полностью автоматизировать процесс получения данных.